%% ============================================================
%%  Section: Introduction  (2026-04-09 v2)
%%  Depositor Discipline in Modern Banking
%%
%%  Structure:
%%    ¶1   Hook + explicit research question
%%    ¶2   Identification challenges + CECL solution
%%    ¶3   Main results: deposit quantities        [was ¶4]
%%    ¶4   Signal validation (large banks)         [was ¶3, compressed to 3 sentences]
%%    ¶5   Main results: financial consequences + sophistication
%%    ¶6   Contribution 1: depositor discipline / market discipline lit
%%    ¶7   Contribution 2: accounting disclosure + CECL consequences lit
%%    ¶8   Contribution 3: deposit insurance design + policy
%%
%%  Changes from intro_section_20260409.tex
%%  (per evaluation .claude/cc/academic-introduction-evaluator/introduction_evaluation_20260408.md):
%%    - ¶1: added explicit "We ask whether..." research question sentence
%%    - ¶3/¶4 swapped: main results now precede signal validation
%%    - ¶3 (formerly ¶4): insured-deposit result clarified to match abstract;
%%      triple-diff coefficient given units and interpretation gloss
%%    - ¶4 (formerly ¶3): compressed from ~9 sentences to 3
%%    - ¶5: terminal misplaced \citep{egan2017deposit, iyer2016tale} removed
%% ============================================================

\section{Introduction}
\label{sec:intro}

Whether uninsured depositors discipline bank risk remains one of the central open
questions in banking.  Theory predicts they should: uninsured depositors bear
full downside risk and have strong incentives to withdraw or demand higher rates
when they detect deterioration in bank financial condition
\citep{calomiris1991role, diamond2001liquidity}.  A large empirical literature
documents that riskier banks pay higher deposit rates and experience larger
uninsured outflows \citep{park1995market, martinez2001depositors,
maechler2006dynamic, iyer2016tale, egan2017deposit, chen2024liquidity}, yet whether these patterns
reflect a causal disciplinary response or equilibrium correlation remains
unsettled.  The question is pressing: against a backdrop of expanding government
protection for uninsured depositors, the Main Street Depositor Protection Act
(S.~2999, reintroduced March 2026) would formally extend coverage to
\$10~million per account, making it important to understand whether this
discipline operates before the incentives that sustain it are further weakened. Historical
evidence suggests the stakes are large, as early state deposit insurance programs in the
United States eliminated market discipline, enabling insured banks to take on excessive
risk \citep{calomiris2019stealing}.
We ask whether uninsured depositors respond to mandatory credit-risk disclosures
by withdrawing balances or demanding higher rates, and we provide causal evidence
on this question using U.S.\ community banks, where institutions below
\$10~billion in assets carry no credible government backstop and identification
conditions are favorable \citep{o1990deposit, flannery1996evidence}.

Achieving clean identification requires resolving three problems that prior work
leaves open.  First, simultaneity: a bank's risk profile, deposit composition,
and depositor behavior co-evolve, so regressions of deposit flows on risk
measures capture equilibrium sorting rather than isolated disciplinary responses.
Second, accounting opacity: under the incurred-loss model, banks could delay
credit loss recognition, systematically understating risk in public disclosures
and limiting depositors' ability to monitor \citep{bushman2012accounting,
bushman2015delayed}.  Third, crisis-period confounding: the samples with the
greatest variation in bank distress are precisely those in which government
guarantees suppress the depositor response \citep{bennett2015market}; financial
stress also raises the liquidity needs of large depositors, making it difficult
to distinguish disciplinary withdrawals from precautionary ones.  We address all
three by exploiting mandatory CECL adoption for community banks in 2023 as a
predetermined information shock.  CECL requires banks to recognize lifetime
expected credit losses on existing loan portfolios at adoption, producing a
one-time adjustment to retained earnings whose magnitude reflects loan vintages
assembled years earlier and therefore cannot be influenced by contemporaneous
managerial decisions or deposit market conditions.  This predetermination is 
central to our identification.  \citet{kim2026current} and \citet{gee2022decision}
confirm that the Day-One adjustment is informative, predicting subsequent
charge-offs and improving analyst forecast accuracy; \citet{nier2006market} and
\citet{chen2022bank} establish that disclosure improvements of this type enable
depositors to act on credit risk signals.

Our main results, estimated on a sample of more than 4,000 community banks over
quarterly observations from 2016 to 2025, establish a clear deposit discipline
pattern.  Banks in the top decile of CECL exposure (adjustment at least 2.5
percent of pre-adoption equity) experienced a 0.394 percentage-point contraction
in uninsured time deposits relative to total assets, roughly 7 percent of the
sample mean.  The continuous treatment confirms the dose-response: each
additional percentage point of CECL exposure reduces uninsured time deposits by
0.072 percentage points of assets.  Insured time deposits, by contrast, show no
statistically distinguishable decline---and if anything trend slightly
upward---consistent with some depositors restructuring balances to stay within
the insurance threshold rather than a uniform funding retreat \citep{martin2026deposit}.  Event studies
confirm that deposit trajectories were parallel across treatment and control banks
for the ten quarters before adoption, and that the divergence emerged at adoption
and persisted through the end of the sample.  A triple-difference specification
stacks the insured and uninsured series in a single panel and adds
bank-by-quarter fixed effects, absorbing all bank-specific time-varying shocks;
the triple interaction of $-1.252$ percentage points of assets per
percentage-point of CECL exposure, which measures the differential decline in
uninsured relative to insured deposits within the same bank and quarter, confirms
a within-bank composition shift rather than a uniform funding retreat.

We confirm that CECL is an information event rather than a mechanical accounting
reclassification using the cohort of large, publicly traded banks that adopted in
2020Q1, the only cohort for which equity prices and credit default swap spreads
are observable around the adoption date.  Equity market capitalization falls
differentially for high-CECL banks beginning in the adoption month, while CDS
spreads widen persistently, reaching approximately 191 basis points per unit of
CECL exposure by month four.  Pre-adoption leads are flat across the full
nine-month window, ruling out anticipation and establishing that CECL is an
information event rather than a mechanical accounting reclassification.

The deposit reallocation carries a real financial cost.  High-CECL banks paid
approximately 8--9 basis points more annually in interest expense relative to
assets, consistent with higher funding costs as they sought to replace outflowing
uninsured balances.  Return on assets fell by roughly 15 basis points per year
and net interest margin compressed by approximately 9 basis points annually,
confirming that higher funding costs were not passed through to borrowers.  The
effects are amplified by depositor sophistication.  Splitting the sample by a
deposit-weighted branch-market index of financial sophistication, constructed
from 2022 FDIC branch data merged to ZIP-code American Community Survey estimates
of educational attainment and IRS investment-income filing rates, we find that
high-CECL banks in high-sophistication markets paid an additional 17--18 basis
points annually in interest expense and suffered an additional 30 basis point
drag on return on assets.  Event studies reveal that the uninsured deposit
decline is concentrated in high-sophistication markets and builds gradually over
the post-adoption quarters.

This paper contributes first to the empirical literature on depositor and market
discipline in banking.  A substantial body of evidence documents that riskier
banks pay higher rates and experience larger uninsured outflows
\citep{park1995market, martinez2001depositors, maechler2006dynamic,
egan2017deposit}, and that market-priced signals contain information about bank
distress \citep{flannery1996evidence, nier2006market}.  The causal
interpretation of these patterns has proven difficult to establish, however.
\citet{bennett2015market} and \citet{berger2015depositors} document that
government guarantees attenuate estimated effects in the crisis episodes that
generate the most variation, and institutions subject to implicit government
support carry systematically smaller estimated discipline effects
\citep{jacewitz2018deposit}.  Unlike prior tests that rely on crisis variation
or cross-sectional risk comparisons in equilibrium, we exploit a predetermined,
mandatory information shock in a non-crisis setting, addressing simultaneity,
accounting opacity, and crisis confounding simultaneously.

This paper also contributes to the accounting and disclosure literature on CECL
and its economic consequences.  \citet{bushman2012accounting} and
\citet{bushman2015delayed} establish that delayed loss recognition under the
incurred-loss model reduced bank transparency and attenuated market discipline;
our findings provide direct evidence that the mandated switch to forward-looking
expected-loss accounting triggered the disciplinary response that opacity had
been suppressing.  Prior work on CECL's real effects focuses on lending behavior
\citep{granja2025current}, provisioning dynamics \citep{kim2026current}, and
decision-usefulness \citep{gee2022decision}.  Unlike these studies, which examine
bank-side responses or analyst outcomes, we trace the full causal chain from
mandatory accounting disclosure to depositor behavior to bank financial
consequences, documenting the funding-market channel through which improved
disclosure standards alter the cost of risk-taking.

Finally, this paper speaks to the debate over deposit insurance design and the
scope of government guarantees.  Prior work establishes that broader insurance
coverage reduces the incentive of uninsured depositors to monitor and exert
pressure on bank risk-taking, weakening the market discipline that complements
regulatory oversight \citep{demirgucc2004market, flannery2019market}.  Our
evidence that uninsured depositors actively monitor and respond to credit risk
information at community banks, even absent a crisis trigger, supports the view
that depositor discipline is a functioning governance mechanism in the
institutional setting where theory predicts it should work best.  Unlike prior
cross-country or crisis-period studies that identify this channel in settings
with substantial confounding from policy interventions and differential
institutional quality, we isolate the mechanism in U.S.\ community banks
operating in normal times.  Policies that move a substantial volume of currently
uninsured deposits inside the insurance perimeter should be evaluated not only
for their run-prevention benefits but also for the monitoring and pricing
discipline they displace.
